%The unsmoothed signal at that index carries a broadband noise excursion from the isolation stage (Section~\ref{sec:galvanic_isolation}) on top of the physiological SCR and would bias the amplitude. EDA

% =============================================================================
% CHAPTER 5 — DISCUSSION
% =============================================================================
%
% Purpose of this chapter (UiO Discussion pattern, seven points):
%   1. What do the results show?
%   2. How do the results fit with other researchers' findings?
%   3. What could the results mean?
%   4. What are the problems and limitations?
%   5. Go back to the big picture from the introduction
%   6. What are the perspectives for future research?
%   7. What are the implications of your findings?
%
% Content policy:
%   Discussion is where interpretation lives. Noise models, mechanism
%   hypotheses, literature comparison, cross-question synthesis, limitations,
%   and implications all belong here, not in Chapter 4.
%
% Structural choice: the seven UiO points map onto seven sections. Sections
% 5.1–5.3 build the interpretive case (results → literature → meaning);
% sections 5.4–5.7 frame what is and isn't known, where it's going, and why
% it matters. The three result threads (isolation, cross-contamination,
% reproducibility/QC) are woven through sections 5.1–5.3 rather than given
% separate subsections, because the interpretive synthesis across threads
% is itself the central contribution.
% =============================================================================

\chapter{Discussion}
\label{ch:discussion}

% -----------------------------------------------------------------------------
% Chapter opening (~half page, before any \section)
% -----------------------------------------------------------------------------
%
% Purpose: establish the interpretive frame. Answers:
%   - What is the central finding in one sentence?
%   - Which threads will be developed, and in which order?
%   - What is the scope of interpretation (bench validation, not clinical)?
%
% Recommended content, two or three paragraphs:
%
% (1) One-sentence headline result. The isolation architecture preserves
%     both signals' morphology and permits simultaneous acquisition without
%     detectable cross-contamination on phasic metrics; a low-frequency EDA
%     coupling is present but does not compromise SCR detection.
%
% (2) Organization preview. Section 5.1 synthesizes what the results show
%     across the three threads; 5.2 places the findings in the literature;
%     5.3 develops the interpretive claims and mechanism hypotheses; 5.4
%     states limitations; 5.5 returns to the gap identified in Chapter 1;
%     5.6 proposes future work; 5.7 states the implications for the
%     biomedical engineering, physiological monitoring, and clinical
%     audiences.
%
% (3) Scope note. Interpretation is bounded by the bench-validation context:
%     five sessions at the palm site with a single subject. Clinical
%     validation at the biceps target site, with a patient cohort, is
%     future work (see 5.6). Statements about reproducibility, mechanism,
%     and implications are made within these limits.
% -----------------------------------------------------------------------------

% -----------------------------------------------------------------------------
% 5.1 What the results show
% -----------------------------------------------------------------------------
%
% This is where the cross-Q synthesis lives. The three-observable ISO224
% convergence belongs here, not in Results. Similarly the Q5/Q6 null
% interpretation and the PSD diagnostic as block-level-not-session-level
% classification.
%
% Three paragraphs recommended:
%   (1) ISO224 characterization — three observables converge on σ_iso ≈ 21 µV
%       RMS (EMG) and low-frequency-concentrated 0.19–0.29 µS (EDA).
%   (2) Cross-contamination — null on primary phasic metrics; block-level
%       EDA coupling at 0.3–1 Hz characterized but does not impair detection.
%   (3) Reproducibility — within-session circuit metrics stable; between-
%       session variability dominated by Session 4. Contamination filter
%       calibration demonstrated at 0.3% exclusion rate.
% -----------------------------------------------------------------------------

\section{What the results show}
\label{sec:disc-synthesis}


% -----------------------------------------------------------------------------
% 5.2 Fit with existing literature
% -----------------------------------------------------------------------------
%
% Requires literature work not yet done. Suggested anchors:
%   - ISO224 datasheet noise specs and any published EMG/EDA front-end noise
%     characterizations using linear isolators.
%   - Prior isolated EMG front-ends (amplitude preservation comparisons).
%   - Multimodal EMG+EDA acquisition systems (if any) and their cross-talk
%     characterizations.
%   - 1/f noise literature for linear isolators (supports the low-frequency
%     spectral shape hypothesis).
%   - EMG contamination-filter conventions (2× within-block median thresholds,
%     or equivalent approaches in the sEMG processing literature).
%
% Approach: do NOT force-fit literature onto every finding. If a finding has
% no direct precedent (e.g., simultaneous EMG+DC-EDA isolation is novel by
% design), state that and cite the closest adjacent work.
% -----------------------------------------------------------------------------

\section{Fit with existing literature}
\label{sec:disc-literature}


% -----------------------------------------------------------------------------
% 5.3 What the results could mean
% -----------------------------------------------------------------------------
%
% Interpretive claims grounded in the data + literature. Candidate content:
%   - ISO224 noise spectrum is low-frequency-concentrated (inferred from EDA
%     amplitude quadrature exceeding σ_local). Physically plausible given 1/f
%     behavior in linear isolators. Discussion-level claim, not proven here.
%   - 0.3–1 Hz EDA coupling: four candidate mechanisms (mechanical coupling
%     through electrode-skin interface, shared power-rail dynamics, cue-
%     triggered physiological arousal, cable-movement artifact). Mechanism
%     cannot be distinguished from current data.
%   - Session 4 asymmetric signature: EMG shows session-wide elevated NF;
%     EDA shows low SCL with strong SCRs. Different physical origins possible
%     (electrode contact on EMG side, subject state on EDA side). Acknowledge
%     limits of retrospective diagnosis.
%   - +0.4 µS SCL shift Combined > EDA-only: either genuine EMG-circuit tonic
%     offset (contamination) or block-order time-on-task effect (confound).
%     Cannot separate at n=3 with fixed block order.
% -----------------------------------------------------------------------------

\section{What the results could mean}
\label{sec:disc-meaning}


% -----------------------------------------------------------------------------
% 5.4 Problems and limitations
% -----------------------------------------------------------------------------
%
% Direct acknowledgment of constraints. Recommended list:
%   - n = 3 for reproducibility (Sessions 3-5). Formal hypothesis testing
%     precluded; between-session SDs are descriptive only.
%   - Block order fixed (EDA-only → EMG-only → Combined). Time-on-task
%     effects inseparable from block effects at n = 3. Directly limits the
%     +0.4 µS SCL shift interpretation.
%   - σ_baseline,session label "circuit noise floor" in methods §3.5.4
%     captures spontaneous EDA activity + slow drift, not pure electronic
%     noise. σ_local is the cleaner circuit-noise metric. σ_session also
%     varies across sessions (S2 = 0.55 vs S3-5 mean ~0.34 µS), limiting
%     cross-session SNR-circuit comparisons; within-session comparisons
%     remain clean.
%   - Ambu wet-gel electrode shelf-life (~1 month after opening); validation
%     data span this window, so electrode quality is not perfectly controlled
%     across sessions.
%   - Palm validation site, not biceps. Clinical motivation is biceps;
%     translation to the target site remains untested.
%   - SCR 100% detection rate is a design outcome (parameters were chosen
%     to achieve it), not an independent metric.
%   - 100/150 Hz "mains harmonic power" values in Q2/Q4 are broadband content
%     in those bands rather than true mains harmonics (filter has no effect
%     on them); reported for completeness, not as interpretable harmonic
%     measurements.
%   - Low-frequency EDA coupling mechanism not identified.
% -----------------------------------------------------------------------------

\section{Problems and limitations}
\label{sec:disc-limitations}


% -----------------------------------------------------------------------------
% 5.5 Back to the introduction: the big picture
% -----------------------------------------------------------------------------
%
% Connect findings to the introduction's gap statement and clinical motivation.
% Key anchors to reference:
%   - Introduction's identified gap: no prior device simultaneously measures
%     EMG and DC EDA from the same muscle site with galvanic isolation.
%   - Clinical motivation: stress/pain monitoring in non-verbal spastic CP
%     patients during passive stretching by physiotherapists.
%   - Collaboration: Dept. of Clinical and Biomedical Engineering (OUS),
%     Dept. of Neurohabilitation (OUS), Dept. of Physics (UiO).
%
% Frame: the gap is now closed at the bench-validation level. Clinical
% validation at the biceps target site remains the next step.
% -----------------------------------------------------------------------------

\section{Back to the big picture}
\label{sec:disc-bigpicture}


% -----------------------------------------------------------------------------
% 5.6 Future research
% -----------------------------------------------------------------------------
%
% Candidate directions ordered by proximity to the thesis contribution:
%   - Clinical validation at biceps with physiotherapy cohort.
%   - Mechanism identification for 0.3–1 Hz EDA coupling: dummy-load tests
%     (distinguish electrical from subject-dependent mechanisms), rail noise
%     measurement (rule in/out NXE2S), cable-routing experiments.
%   - Extended ISO224 spectral characterization (direct measurement to
%     confirm low-frequency concentration hypothesis).
%   - Larger reproducibility cohort (n > 3) with randomized block order to
%     enable formal statistical testing and decouple block effects from
%     time-on-task.
%   - Alternative electrode technologies (dry electrodes, extended-shelf-
%     life wet gel) for longer-duration clinical deployments.
%   - Co-located accelerometry for motion-artifact rejection if mechanical
%     coupling is confirmed as the 0.3–1 Hz mechanism.
% -----------------------------------------------------------------------------

\section{Perspectives for future research}
\label{sec:disc-future}


% -----------------------------------------------------------------------------
% 5.7 Implications
% -----------------------------------------------------------------------------
%
% Three audiences:
%   - Biomedical front-end designers: ISO224 is a viable component for
%     simultaneous physiological signal acquisition, with quantified noise
%     cost (~21 µV RMS EMG, ~0.2 µS RMS EDA low-frequency-concentrated).
%   - Stress/pain monitoring in non-verbal patients: combined EMG+EDA from
%     the same muscle site provides two independent arousal signals without
%     mutual interference on phasic metrics — meeting the key methodological
%     requirement stated in the introduction.
%   - OUS Neurohabilitation clinical team: the device is bench-validated
%     and ready for subject testing at the biceps target site.
% -----------------------------------------------------------------------------

\section{Implications}
\label{sec:disc-implications}

\begin{comment}
    
\begin{itemize}
    \item What do the results show?
    \item How do the results fit in with other researchers findings?
    \item What could the results mean?
    \item What are the problems and limitations
    \item Go back to the big picture from the introduction
    \item What are the perspectives for future research?
    \item What are the implications of your findings?
\end{itemize}
\end{comment}